\documentclass[a4paper,11pt]{article}
\usepackage[utf8]{inputenc}
\usepackage[T1]{fontenc}
\usepackage[french]{babel}
\usepackage{lmodern}
\usepackage{amsmath,amssymb,amsthm}
\usepackage{mathrsfs}
\usepackage{enumitem}
\usepackage{multicol}

\setlist[enumerate]{label=\textbf{\textcolor{Couleur8}{\arabic*.}}, left=1em}

% Si vous voulez aussi modifier les listes enumerate imbriquées :
\setlist[enumerate,2]{label=\textcolor{Couleur8}{\alph*)}}
\setlist[enumerate,3]{label=\textcolor{Couleur8}{\roman*)}}
\setlist[enumerate,4]{label=\textcolor{Couleur8}{\Alph*)}}
\usepackage{graphicx}
\usepackage{xcolor}
\usepackage{tcolorbox}
\usepackage{geometry}
\usepackage{fancyhdr}
\usepackage{titlesec}
\usepackage{cancel}
\usepackage{ifthen}
\usepackage{tikz}
\usetikzlibrary{arrows.meta}
\usepackage{tkz-tab}
\usepackage{eurosym}

% OPTION PROF/ÉLÈVE
% Décommentez UNE des deux lignes suivantes :
\newboolean{prof}\setboolean{prof}{true}   % Version professeur (avec corrections des devoirs)
\newboolean{prof}\setboolean{prof}{true}  % Version élève (sans corrections des devoirs)

\definecolor{Couleur1}{HTML}{4A5D7A} %Th
\definecolor{Couleur2}{HTML}{A5B5D6} %Prop
\definecolor{Couleur3}{HTML}{A8C68A} %Méthode
\definecolor{Couleur6}{HTML}{8A9CC1} %Def
\definecolor{Couleur7}{HTML}{E6B459} %Rq Voc
\definecolor{Couleur8}{HTML}{D36740} %Titres
\definecolor{correction}{HTML}{D36740} % Couleur pour les corrections
\definecolor{coopmathscorr}{HTML}{F15929} % Couleur corrections coopmaths

% Configuration des marges
\geometry{
	top=2cm,
	bottom=2cm,
	inner=1.5cm,
	outer=1.5cm,
}

% Police
\renewcommand{\familydefault}{\sfdefault}

% Compteurs
\newcounter{seancenum}
\newcounter{seancenumD}
\setcounter{seancenum}{1}
\setcounter{seancenumD}{1}

% Configuration des en-têtes et pieds de page
\pagestyle{fancy}
\fancyhf{}
\ifthenelse{\boolean{prof}}{
    \fancyhead[L]{\textcolor{Couleur8}{\textbf{Corrections Automatismes - Classe de seconde - Version Professeur}}}
}{
    \fancyhead[L]{\textcolor{Couleur8}{\textbf{Corrections Devoirs - Séance 24 - Classe de seconde}}}
}
\fancyhead[R]{\textcolor{Couleur8}{\textbf{2025-2026}}}
\fancyfoot[L]{\textcolor{Couleur8}{Mme Bertail et Mme Pinto}}
\fancyfoot[R]{\textcolor{Couleur8}{\thepage}}
\renewcommand{\headrulewidth}{1pt}
\renewcommand{\headrule}{\hbox to\headwidth{\color{Couleur8}\leaders\hrule height \headrulewidth\hfill}}

% Environnements personnalisés
\newtcolorbox{seance}[1]{
    colback=Couleur6!10,
    colframe=Couleur1!65,
    fonttitle=\bfseries\large,
    title=Correction Séance \theseancenum #1,
    left=5mm,
    right=5mm,
    top=3mm,
    bottom=3mm,
    arc=0mm,
    after=\stepcounter{seancenum}
}

\newtcolorbox{seanceD}[1]{
    colback=Couleur6!5,
    colframe=Couleur6!45,
    fonttitle=\bfseries\large,
    title=Correction Devoirs - Séance \theseancenumD #1,
    left=5mm,
    right=5mm,
    top=3mm,
    bottom=3mm,
    arc=0mm,
    after=\stepcounter{seancenumD}
}

\begin{document}

\setcounter{seancenumD}{24}
\newpage
\begin{seanceD}{} %24

\begin{enumerate}
\item On procède par étapes successives :\\
      On commence par isoler $-4x$ dans le membre de gauche en ajoutant
      $-25$ dans chacun des membres, puis on divise
      par $-4$ pour obtenir la solution : \\
       $\begin{aligned}
       -4x+25&=-7\\
      -4x&=-7-25\\
      -4x&=-32\\
      x&=\dfrac{-32}{-4}\\
      x&=8
      \end{aligned}$\\
      La solution de l'équation est : ${\color[HTML]{f15929}\boldsymbol{8}}$.
      
	\item $ \dfrac{25}{35}=\dfrac{5\times 5}{7\times 5}
      ={\color[HTML]{f15929}\boldsymbol{\dfrac{5}{7}}}$
	\item La notation scientifique est de la forme $a\times 10^{n}$ avec $1\leqslant a <10$ et $n$ un entier relatif.\\
              Ici : $0{,}002\,53=\underbrace{{\color[HTML]{f15929}\boldsymbol{2{,}53}}}_{1\leqslant 2{,}53 <10}{\color[HTML]{f15929}\boldsymbol{\times}} {\color[HTML]{f15929}\boldsymbol{10^{-3}}}$. 
	\item On utilise la formule $a^n\times a^m=a^{n+m}$ avec $a=2$, $n=4$ et $p=2$.\\
            $2^{4}\times 2^{2}=2^{4+2}=2^{{\color[HTML]{f15929}\boldsymbol{6}}}$
	\item Pour $x=-3$, on obtient : 
            $3+4x=3+4\times(-3)={\color[HTML]{f15929}\boldsymbol{-9}}$.
            
	\item Une unité correspond à $4$ carreaux, la ligne brisée mesure $7$ carreaux, soit $\dfrac{{\color[HTML]{f15929}\boldsymbol{7}}}{{\color[HTML]{f15929}\boldsymbol{4}}}$ u.l. 
	\item Comme il y a $1$ multiple de $5$, la probabilité d'obtenir un multiple de $5$ est ${\color[HTML]{f15929}\boldsymbol{\dfrac{1}{6}}}$.
	\item L'équation de la droite est de la forme $y=mx+p$.\\
       Le coefficient directeur est $m$ et l'ordonnée à l'origine est $p$. \\
      Le coefficient directeur de la droite est donc $m={\color[HTML]{f15929}\boldsymbol{-3}}$.
	\item Les coordonnées du milieu sont données par la moyenne des abscisses et la moyenne des ordonnées : \\
      $x_M=\dfrac{4+6}{2}={\color[HTML]{f15929}\boldsymbol{5}}$ et $y_M=\dfrac{4+8}{2}={\color[HTML]{f15929}\boldsymbol{6}}$.\\
      Ainsi,  $M({\color[HTML]{f15929}\boldsymbol{5\,;\,6}})$.
	\item On utilise l'égalité remarquable $(a+b)^2=a^2+2ab+b^2$ avec $a=x$ et $b=4$.\\
      $(x+4)^2=x^2+2 \times 4 \times x+4^2={\color[HTML]{f15929}\boldsymbol{x^2+8x+16}}$
	\item On utilise l'égalité remarquable $a^2-b^2=(a+b)(a-b)$ avec $a=x$ et $b=4$.\\
      $x^2-16=x^2-4^2={\color[HTML]{f15929}\boldsymbol{(x-4)(x+4)}}$
	\item $10$ viennent au lycée à vélo, donc $13$ ne viennent pas au lycée à vélo.\\
      La proportion d’élèves de cette classe qui ne viennent pas à vélo est donc ${\color[HTML]{f15929}\boldsymbol{\dfrac{13}{23}}}$.
	\item Le nombre de solution de l'équation $f(x)=4$ est le nombre de points d'intersection entre la droite d'équation $y=4$ et la courbe de $f$. \\La droite horizontale d'équation $y=4$ n'a aucun point d'intersection avec la courbe de $f$. \\
          Ainsi l'équation $f(x)=4$ a ${\color[HTML]{f15929}\boldsymbol{0}}$ solution.
	\item Pour lire l'image de $3$, on place la valeur de $3$ sur l'axe des abscisses (axe de lecture  des antécédents) et on lit
           son image  sur l'axe des ordonnées (axe de lecture des images). \\
           On obtient :  $f(3)={\color[HTML]{f15929}\boldsymbol{1}}$.
	\item La fonction est positive ou nulle lorsque les images sont positives ou nulles.\\
        Graphiquement, les images sont positives ou nulles  lorsque la courbe se situe sur ou au-dessus  de l'axe des abscisses, soit sur l'intervalle  
        ${\color[HTML]{f15929}\boldsymbol{[0\,;\,3]}}$
        
\end{enumerate}


\end{seanceD}

\end{document}